\section{Discussion}
\label{sec:discussion}

\subsection{Why Observation Design Matters}

The O0--O1 factorial separates observation design from estimator complexity. With the RLS kernel unchanged, O1 reduced capacitance MAPE from $7.0366\%$ to $0.3727\%$, ESR MAPE from $1.1970\%$ to $0.2233\%$, and mean internal convergence from 665.41 to 13.21 cycles. This attribution agrees with the circuit physics: the switching edge exposes the ESR drop, while the in-state charge balance exposes inverse capacitance. More elaborate recursion cannot reliably compensate for observations that mix these effects or compare mismatched timestamps.

The scoped literature matrix contains no disclosed method combining the \Cuk{} transfer-capacitor target, $+i_{L1}\leftrightarrow-i_{L2}$ commutation, timestamp-reconstructed edge information, and direction-specific charge/edge updates. This is a bounded review finding, not a universal first-in-literature claim.

\subsection{Why \TSRLS{} Is Primary}

\TSRLS{} is primary because it combines low ESR error, rapid convergence, low operation count, and the better aggregate ramp-tracking tradeoff. This choice neither makes RLS a novelty nor implies superiority in every metric. \TSKF{} remains valuable for confidence-aware reporting: its joint coverage is much closer to nominal, but its slower ramp response and retained abrupt-step failures limit it as the sole point estimator. An adaptive change detector or process-noise policy could alter that tradeoff but lies outside the evaluated design.

\subsection{Limitations}

The evidence has six principal boundaries.
\begin{enumerate}
\item Identifiability and updates require valid CCM; DCM and weak-excitation intervals hold the estimate. The method assumes availability of both inductor currents; applications lacking either measurement require an additional sensor or independently validated current reconstruction.
\item The reported ESR is an effective condition indicator at the modeled switching, thermal, and frequency conditions, not a temperature-independent aging variable.
\item The $0.1$--$100$~s trajectories are synthetic or trace-derived bandwidth tests, not run-to-failure data; abrupt discontinuities are recovery stress tests rather than aging events.
\item Hardware validation remains open for edge ringing, isolation/probing, AFE common-mode rejection, and production calibration.
\item The F28379D study is device-realistic simulation; target compilation and measured WCET have not been demonstrated.
\item The auxiliary \TSRLS{} intervals are under-covered, while \TSKF{} retains the reported dynamic limitations; confidence calibration and point-tracking performance therefore remain a tradeoff.
\end{enumerate}

\subsection{Experimental Roadmap}

The next stage should use a controlled C/ESR network in the energy-transfer branch, the specified F28379D schedule, separate absolute- and edge-voltage paths, and synchronized inductor-current measurements. Tests should sweep input voltage, duty ratio, load, and temperature; characterize ringing and AFE rejection; verify calibration against offline impedance measurements; and measure WCET before firmware release.
